\chapter{Những Điều Thầy Giáo Cũng Đang Học}
\markboth{Những Điều Thầy Giáo Cũng Đang Học}{What Teachers Are Still Learning Too}

\section{Không Phải Lúc Nào Cũng Đúng}
\begin{truyen}{Không Phải Lúc Nào Cũng Đúng}{Not Being Right All the Time}
\chuhoa{D}{ạy lâu mới thấy có những lần mình tưởng mình hiểu học sinh rồi, nhưng hóa ra mình mới chỉ hiểu một lát cắt.}
\textit{(Teaching long enough reveals how often we think we understand a student when we have only seen one slice.)}

Có lúc thầy cô đoán sai, nóng sai, hoặc đánh giá thiếu.
\textit{(Teachers sometimes guess wrong, react wrongly, or judge incompletely.)}

Biết mình không phải lúc nào cũng đúng không làm nghề dạy yếu đi.
\textit{(Knowing that we are not always right does not weaken teaching.)}

Nó làm người dạy bớt cứng và nhìn học trò thật hơn.
\textit{(It makes teachers less rigid and more truthful in seeing students.)}

\end{truyen}

\begin{baihoc}
Sự khiêm tốn của người dạy cũng là một phần của bài học.
\textit{(A teacher’s humility is also part of the lesson.)}
\end{baihoc}

\begin{ghinhoanh}
\textbf{Short English to remember:}

A teacher’s humility is also part of the lesson.

\end{ghinhoanh}

\begin{tuhocnhanh}
\textbf{Gợi ý để dạy và tự nghĩ / Teaching and reflection prompts}

1. Vì sao người dạy cũng cần chấp nhận mình có thể sai? \textit{Why must teachers accept that they can be wrong?}

2. Điều này giúp lớp học thế nào? \textit{How does this help the classroom?}

3. Mình có đang đủ khiêm tốn trước học sinh không? \textit{Are we humble enough before students?}
\end{tuhocnhanh}
\ngancach

\section{Nóng Giận Không Phải Lúc Nào Cũng Là Uy}
\begin{truyen}{Nóng Giận Không Phải Lúc Nào Cũng Là Uy}{Anger Is Not Always Authority}
\chuhoa{C}{ó lúc người dạy lớn tiếng và lớp im ngay. Cảm giác đó rất dễ làm ta tưởng mình vừa lập lại trật tự đúng cách.}
\textit{(Sometimes a teacher raises the voice and the class falls silent immediately. That feeling can tempt us to believe order has been restored in the right way.)}

Nhưng sự im do sợ khác rất xa sự im để học.
\textit{(But silence created by fear is very different from silence that makes learning possible.)}

Uy tín thật không chỉ làm học sinh nín. Nó làm học sinh tin.
\textit{(Real authority does not only silence students. It helps them trust.)}

Có những lần người dạy phải học cách hạ giọng để nghề mình cao hơn.
\textit{(There are times teachers must learn to lower the voice so that the profession itself stands higher.)}

\end{truyen}

\begin{baihoc}
Không phải sự im nào trong lớp cũng là dấu hiệu tốt.
\textit{(Not every quiet classroom is a good sign.)}
\end{baihoc}

\begin{ghinhoanh}
\textbf{Short English to remember:}

Not every quiet classroom is a good sign.

\end{ghinhoanh}

\begin{tuhocnhanh}
\textbf{Gợi ý để dạy và tự nghĩ / Teaching and reflection prompts}

1. Sự im này đến từ học hay từ sợ? \textit{Is this silence coming from learning or from fear?}

2. Uy tín thật khác áp lực ở đâu? \textit{How is real authority different from pressure?}

3. Mình có thể điều chỉnh điều gì trong cách phản ứng? \textit{What can we adjust in our responses?}
\end{tuhocnhanh}
\ngancach

\section{Mỗi Em Cần Một Cánh Cửa Khác}
\begin{truyen}{Mỗi Em Cần Một Cánh Cửa Khác}{Each Student Needs a Different Door}
\chuhoa{C}{ó em chỉ cần một lời khen là sáng ra. Có em phải cần thời gian. Có em phải cần một lần được tin trước khi chịu mở lòng.}
\textit{(Some students brighten with one word of praise. Some need time. Some need to be trusted once before opening up.)}

Dạy lâu mới thấy không có một chìa khóa cho tất cả.
\textit{(Teaching long enough shows that there is no single key for everyone.)}

Muốn chạm được học trò, người dạy phải chịu khó thử nhiều cửa.
\textit{(To reach students, teachers must be willing to try many doors.)}

Nghề này mệt ở chỗ đó, nhưng cũng đẹp ở chỗ đó.
\textit{(That is one exhausting part of this work, and also one of its beauties.)}

\end{truyen}

\begin{baihoc}
Dạy học không chỉ là truyền đạt. Còn là đi tìm đúng cánh cửa để một đứa trẻ chịu bước ra.
\textit{(Teaching is not only delivering knowledge. It is also finding the right door through which a child is willing to step out.)}
\end{baihoc}

\begin{ghinhoanh}
\textbf{Short English to remember:}

Teaching is also about finding the right door for each child.

\end{ghinhoanh}

\begin{tuhocnhanh}
\textbf{Gợi ý để dạy và tự nghĩ / Teaching and reflection prompts}

1. Vì sao không thể dùng một cách cho mọi học sinh? \textit{Why can’t one method fit every student?}

2. “Cánh cửa” của học sinh có thể là gì? \textit{What might a student’s “door” be?}

3. Mình có đủ kiên nhẫn để đi tìm không? \textit{Are we patient enough to keep looking?}
\end{tuhocnhanh}
\ngancach

\section{Dạy Học Cũng Là Học Lại Làm Người}
\begin{truyen}{Dạy Học Cũng Là Học Lại Làm Người}{Teaching Is Also Relearning How to Be Human}
\chuhoa{S}{au nhiều năm đứng lớp, có người nhận ra mình không chỉ đang dạy học sinh. Mình cũng đang được học sinh dạy lại rất nhiều.}
\textit{(After many years in the classroom, some teachers realize they are not only teaching students. Students are also teaching them many things.)}

Dạy lại cho người lớn về kiên nhẫn, về sự vụng về, về lòng tự trọng bé nhỏ, và về những vết thương không nói ra.
\textit{(Students reteach adults patience, clumsiness, small forms of dignity, and wounds that are rarely spoken.)}

Nếu nhìn đủ sâu, nghề dạy không chỉ cho bài giảng. Nó cho một người lớn cơ hội bớt cứng và tử tế hơn.
\textit{(Seen deeply enough, teaching does not only produce lessons. It gives an adult the chance to become less rigid and more humane.)}

Có lẽ vì thế mà nghề này mệt, nhưng khó bỏ.
\textit{(Perhaps that is why this profession is exhausting, yet hard to leave.)}

\end{truyen}

\begin{baihoc}
Người thầy tốt không chỉ dạy kiến thức. Họ cũng để nghề dạy sửa mình từng ngày.
\textit{(A good teacher does not only teach knowledge. They also let the profession refine them day by day.)}
\end{baihoc}

\begin{ghinhoanh}
\textbf{Short English to remember:}

A good teacher lets the profession refine them day by day.

\end{ghinhoanh}

\begin{tuhocnhanh}
\textbf{Gợi ý để dạy và tự nghĩ / Teaching and reflection prompts}

1. Học sinh đã dạy ngược lại điều gì cho người lớn? \textit{What do students teach adults in return?}

2. Vì sao nghề dạy học vừa mệt vừa khó bỏ? \textit{Why is teaching exhausting yet hard to leave?}

3. Mình đang để nghề sửa mình ở điểm nào? \textit{Where is the profession refining us right now?}
\end{tuhocnhanh}
\ngancach

\section{Không Phải Bài Nào Mình Dạy Cũng Tới Được Hết}
\begin{truyen}{Không Phải Bài Nào Mình Dạy Cũng Tới Được Hết}{Not Every Lesson Reaches Everyone}
\chuhoa{N}{gười mới dạy thường nghĩ nếu chuẩn bị đủ kỹ, nói đủ rõ, làm đủ tâm thì bài sẽ chạm được hầu hết học sinh. Dạy lâu mới hiểu có những hôm bài vẫn hụt.}
\textit{(New teachers often believe that if they prepare well, explain clearly, and care deeply, the lesson will reach most students. Experience teaches that some days it still misses.)}

Không phải lúc nào hụt cũng là lỗi do mình thiếu cố gắng.
\textit{(A missed lesson is not always caused by a lack of effort.)}

Lớp học có quá nhiều thứ đi vào bài giảng ngoài kiến thức: tâm trạng, hoàn cảnh, sức khỏe, chuyện gia đình, và thời điểm.
\textit{(Too many things enter a lesson beyond knowledge: mood, circumstance, health, family life, and timing.)}

Biết vậy để bớt kiêu, và cũng để bớt tự trách vô lý.
\textit{(Knowing this helps us become less proud and less unfairly self-blaming.)}

\end{truyen}

\begin{baihoc}
Nghề dạy cần sự tận tâm, nhưng cũng cần sự thật thà trước giới hạn của mình.
\textit{(Teaching requires devotion, but also honesty about one’s limits.)}
\end{baihoc}

\begin{ghinhoanh}
\textbf{Short English to remember:}

Teaching requires devotion, but also honesty about one’s limits.

\end{ghinhoanh}

\begin{tuhocnhanh}
\textbf{Gợi ý để dạy và tự nghĩ / Teaching and reflection prompts}

1. Vì sao có bài dạy vẫn hụt? \textit{Why do some lessons still miss?}

2. Điều gì ngoài kiến thức đi vào một tiết học? \textit{What enters a lesson besides knowledge?}

3. Mình có đang tự trách quá mức không? \textit{Am I blaming myself too much?}
\end{tuhocnhanh}
\ngancach

\section{Có Lúc Mình Nói Nhanh Hơn Mình Hiểu}
\begin{truyen}{Có Lúc Mình Nói Nhanh Hơn Mình Hiểu}{Sometimes We Speak Faster Than We Understand}
\chuhoa{N}{ghề dạy tạo cho người ta thói quen phản ứng nhanh. Nhưng có lúc sự nhanh đó làm mình nói ra trước khi hiểu đủ học sinh đang cần gì.}
\textit{(Teaching builds a habit of responding quickly. But sometimes that speed makes us speak before we understand what a student truly needs.)}

Có khi học sinh cần giải thích. Có khi cần bị dừng lại. Có khi chỉ cần được nghe đến hết.
\textit{(Sometimes a student needs explanation. Sometimes they need to be stopped. Sometimes they only need to be heard to the end.)}

Người dạy trưởng thành không phải là người lúc nào cũng có câu trả lời nhanh.
\textit{(A mature teacher is not one who always has a fast answer.)}

Mà là người biết chậm lại cho đúng chỗ.
\textit{(It is someone who knows when to slow down in the right place.)}

\end{truyen}

\begin{baihoc}
Chậm một chút để hiểu kỹ có khi cứu được một cuộc nói chuyện khỏi đi sai hướng.
\textit{(Slowing down to understand can save a conversation from moving in the wrong direction.)}
\end{baihoc}

\begin{ghinhoanh}
\textbf{Short English to remember:}

Slowing down to understand can save a conversation.

\end{ghinhoanh}

\begin{tuhocnhanh}
\textbf{Gợi ý để dạy và tự nghĩ / Teaching and reflection prompts}

1. Vì sao người dạy dễ nói nhanh? \textit{Why do teachers tend to respond quickly?}

2. Chậm lại ở đâu là cần thiết? \textit{Where is slowing down necessary?}

3. Mình có đang nói nhanh hơn hiểu không? \textit{Am I speaking faster than I understand?}
\end{tuhocnhanh}
\ngancach

\section{Có Những Em Làm Mình Sửa Cả Cách Dạy}
\begin{truyen}{Có Những Em Làm Mình Sửa Cả Cách Dạy}{Some Students Make Us Rethink Our Entire Teaching Style}
\chuhoa{C}{ó em học không vào bằng cách cũ, nghe không được kiểu cũ, và phản ứng rất khác với số đông. Ban đầu người dạy dễ nghĩ em khó. Dạy lâu mới hiểu có khi cách mình đang dùng mới là thứ chưa hợp.}
\textit{(Some students do not learn through the old way, do not respond to the usual tone, and react differently from the majority. At first the teacher may think the student is difficult. With time, we realize the method may be the real problem.)}

Học sinh không chỉ học từ thầy cô. Các em cũng ép người dạy phải học lại nghề của mình.
\textit{(Students do not only learn from teachers. They also force teachers to relearn their profession.)}

Đó là phần vất vả, nhưng cũng là phần giúp nghề này không chết cứng.
\textit{(That is the tiring part, but also the part that keeps the profession from becoming rigid.)}

Có những em đến để học. Có những em đến để dạy lại mình cách dạy.
\textit{(Some students come to learn. Some come to reteach us how to teach.)}

\end{truyen}

\begin{baihoc}
Người dạy lớn lên thật sự khi dám sửa cả chính cách mình từng tin là đúng.
\textit{(Teachers truly grow when they dare to revise even the methods they once believed were right.)}
\end{baihoc}

\begin{ghinhoanh}
\textbf{Short English to remember:}

Teachers truly grow when they dare to revise even their trusted methods.

\end{ghinhoanh}

\begin{tuhocnhanh}
\textbf{Gợi ý để dạy và tự nghĩ / Teaching and reflection prompts}

1. Học sinh nào từng làm mình phải đổi cách dạy? \textit{Which students have forced me to change my method?}

2. Vì sao sự thay đổi đó quan trọng? \textit{Why does that change matter?}

3. Mình có đủ dẻo để học lại nghề không? \textit{Am I flexible enough to relearn my profession?}
\end{tuhocnhanh}
\ngancach

\section{Người Dạy Cũng Cần Được Nói “Anh Đã Làm Được Rồi”}
\begin{truyen}{Người Dạy Cũng Cần Được Nói “Anh Đã Làm Được Rồi”}{Teachers Also Need to Hear “You Did Well”}
\chuhoa{N}{ghề dạy quen với việc động viên học sinh, nhưng người dạy cũng có lúc rất cần một câu công nhận.}
\textit{(Teachers are used to encouraging students, but teachers too sometimes need one sentence of recognition.)}

Không phải vì yếu đuối. Mà vì ai làm việc với con người lâu cũng có lúc mòn.
\textit{(Not because they are weak, but because anyone working with people for a long time wears down at times.)}

Một lời nói đúng lúc từ đồng nghiệp, phụ huynh, hay học sinh cũ có thể làm mình đứng lại vững hơn rất nhiều.
\textit{(A timely word from a colleague, parent, or former student can steady a teacher greatly.)}

Người dạy cũng là con người. Điều đó càng nhớ rõ thì nghề này càng bền.
\textit{(Teachers are human too. The more clearly we remember that, the more sustainable the profession becomes.)}

\end{truyen}

\begin{baihoc}
Muốn người dạy tiếp tục cho đi lành mạnh, cũng cần có cách nuôi lại họ.
\textit{(If we want teachers to keep giving in healthy ways, we also need ways to replenish them.)}
\end{baihoc}

\begin{ghinhoanh}
\textbf{Short English to remember:}

If we want teachers to keep giving well, we must also replenish them.

\end{ghinhoanh}

\begin{tuhocnhanh}
\textbf{Gợi ý để dạy và tự nghĩ / Teaching and reflection prompts}

1. Vì sao người dạy cũng cần được công nhận? \textit{Why do teachers also need recognition?}

2. Một lời công nhận đúng lúc có thể làm gì? \textit{What can timely recognition do?}

3. Mình có đang biết cách nuôi lại chính mình không? \textit{Do I know how to replenish myself?}
\end{tuhocnhanh}
